\subsection{Nintendo}

\subsubsection{Official Policy}

Nintendo's \textit{Game Content Guidelines for Online Video
\& Image Sharing Platforms}, updated September 2, 2024,
permit qualifying gameplay videos and livestreams.

The guidelines encourage original creative input or
commentary and recognize platforms including, but not
limited to, YouTube and Twitch.

Official policy:
\url{https://www.nintendo.co.jp/networkservice_guideline/en/index.html}

\subsubsection{Applicability to Silima}

We believe Silima is consistent with the broadcasting
activities contemplated by Nintendo's guidelines for
the following reasons:

\begin{itemize}
    \item \textbf{Broadcasting Platform:}
    Silima is designed to transmit gameplay from a host
    computer to connected viewers, functioning similarly
    to established livestreaming platforms.

    \item \textbf{Creative Contribution:}
    Gameplay is collectively controlled by participants
    through voting and shared inputs, producing an
    unpredictable experience shaped by the audience.

    \item \textbf{Interactive Commentary:}
    Participants can react to gameplay through emojis
    and predefined messages shared with other viewers,
    adding contextual commentary and social interaction.

    \item \textbf{Noncommercial Use:}
    Silima does not charge participants for gameplay
    or monetize broadcasts.
\end{itemize}

We believe these characteristics provide the creative
and interactive contributions contemplated by Nintendo's
guidelines, rather than merely reproducing unmodified
gameplay footage.

However, Nintendo expressly limits its guidelines to
individual consumers and restricts certain uses involving
unauthorized game software. The guidelines do not
explicitly authorize institutional broadcasting through
Silima, and compliance with these additional restrictions
must be evaluated separately.
